% 来源：OEBPS/Text/01_42.xhtml（OPF spine 顺序；仅第一册）
\cleardoublepage
\BookEntry[毛姆《人生的枷锁》]{毛姆《人生的枷锁》}

\BookCover{images/书单26人生的枷锁.jpg}{[英] 毛姆著| 张柏然等译| 上海译文出版社2007年版}

\begin{BookSummaryBox}
\textbf{毛姆具有自传色彩的代表性作品}
\par
\end{BookSummaryBox}

\subsubsection*{这本书谈什么？}

小说《人生的枷锁》中，主人公菲利普·凯里自幼父母双亡，不幸又先天残疾，他在冷漠而陌生的环境中度过了童年，性格孤僻又敏感。在寄宿学校度过的岁月让他饱受了不合理的学校制度的摧残，而当他走入社会后，又承受了爱情苦痛的折磨。在坎坷的人生道路上，他每跨一步，都要付出艰辛的挣扎，但思想和个性都独立不羁的凯里，一直努力挣脱宗教和小市民意识这两条禁锢自己精神的桎梏，力图在混沌纷扰的生活旋涡中，寻求人生的真谛。
\par

\subsubsection*{芒格为什么谈起这本书？}

芒格在南加州大学GOULD法学院毕业典礼的演讲中，谈道：
\par

\begin{quotation}
正确的爱应该以仰慕为基础，而且我们应该去爱那些对我们有教育意义的先贤。萨默赛特·毛姆在他的小说《人生的枷锁》中描绘的爱是一种有病的爱。那是一种病，如果你们发现自己有那种病，应该赶快把它治好。
\par
\end{quotation}

1915年，威廉·萨默赛特·毛姆出版了自传体小说《人生的枷锁》，这本书通常被认为是他最出色的作品。小说的主角菲利普遇到了伦敦的女服务员米尔德莱德。米尔德莱德瞧不起菲利普，但菲利普疯狂地爱上了她。菲利普知道他自己很蠢，也很讨厌自己。他把所有钱都给了米尔德莱德，而得到的回报却是憎恶和羞辱。毛姆如此描绘这段关系：
\par

\begin{quotation}
爱情就如他腿上的寄生虫，吮吸着他生命的热血，维持其可恶的生存；它拼命地吸取他的精力，使他对别的一切都提不起兴趣。
\par
\end{quotation}

\subsubsection*{《人生的枷锁》作者毛姆如是说}

“世界上最大的折磨也莫过于在爱的同时又带着藐视了。”
\par

“他全部生活曾经是亦步亦趋地跟在别人的理想后面，那是他们用语言和书本灌输给他的思想，而绝非是他自己心里的愿望。过去左右着他生活道路的那些他认为是应该做的，而不是他全心全意希望去做的事情。现在他不耐烦地把这所有一切都抛弃了……也许对幸福的俯就是自认失败，但是，这种失败比许多次胜利都要强。”
\par

“一个人要意识到自我的存在，也非得经历同样的痛苦不可。在此，差别在于，虽然每个人同样认识到自己的身体是一个独立而完整的有机体，但是，并不是每个人都同样认识到自己是一个完整、独立的人的存在。这种离群索居的感觉在青春期尤为明显。可是这种感觉，并没有发展到使个人和同伴之间的差别达到今人一目了然的明显程度。只有像蜂巢里的蜜蜂那样很少自我意识的人，才是生活的幸运儿，因为他们最有机会获得幸福。他们集体行动，群起群居，而他们的欢乐也只因为大家共享才成其欢乐。”
\par
